\subsection{معرفی و انگیزه}
این مکانیزم بر مبنای مبانی نظری پیشرفته در ادبیات توجه و فشرده‌سازی \cite{kim2022learned} توسعه یافته است.

کاهش ابعاد و حجم محاسباتی شبکه‌های عصبی عمیق بدون افت بازنمایی معنایی، یکی از چالش‌های بنیادین در استقرار مدل‌های ترنسفورمر روی سخت‌افزارهای با حافظه محدود است. روش‌های کلاسیک هرس ساختارنیافته هرچند می‌توانند افزونگی پارامتری را تا حد زیادی حذف کنند، اما به دلیل عدم انطباق با واحدهای پردازش تانسوری مدرن (\lr{Tensor Cores})، عملاً به شتاب‌بخشی محاسباتی و کاهش زمان تاخیر منجر نمی‌شوند. از سوی دیگر، هرس متراکم کم‌مرتبه (\lr{Low-Rank Decomposition}) اطلاعات فضایی کانال‌ها را دگرگون ساخته و در نسبت‌های فشرده‌سازی بالا با افت دقت مواجه می‌شود.

برای حل این چالش، چارچوب فشرده‌سازی \lr{ARMOR: Adaptive Resilient Matrix Orthogonal Rank Pruning} معرفی شده است. این موتور فشرده‌سازی، ساختار ماتریس وزن‌ها را به صورت ترکیبی نامتقارن از یک هسته تنک با ساختار سخت‌افزاری \lr{2:4} (\lr{2:4 Structured Sparsity}) و یک تصحیح‌کننده باقیمانده کم‌مرتبه (\lr{Low-Rank Residual Correction}) مدل‌سازی می‌کند. همچنین این ماژول با استفاده از معیار حساسیت تیلور انعطاف‌پذیر (\lr{Resilient Taylor Saliency}) با ساختار پشتیبان سه‌سطحی، تضمین می‌کند که مهم‌ترین عناصر وزن تحت هر شرایط داده‌ای شناسایی شده و اختلال کمینه‌ای به سیگنال‌های بازنمایی مدل وارد گردد.
